% 02_projective_qsp.tex
\section{Quantum signal processing on the projective line}
\label{sec:projective}

\subsection{The qubit as a point of $\CP$}
\label{subsec:encoding}

A single-qubit pure state, taken up to global phase, is a point of
the complex projective line $\CP=\C\cup\{\infty\}$.  We fix the
identification by stereographic projection from the north pole of
the Bloch sphere,
\begin{equation}
	\mathcal{S}\colon\ket{\psi}\ \longmapsto\
	\frac{\bra{\psi}\PauliX\ket{\psi}-i\bra{\psi}\PauliY\ket{\psi}}
	{1-\bra{\psi}\PauliZ\ket{\psi}}\ \in\ \CP,
	\label{eq:decode}
\end{equation}
with inverse
\begin{equation}
	\ket{z}=\frac{1}{\sqrt{1+|z|^2}}\bigl(z\ket{0}+\ket{1}\bigr).
	\label{eq:encode}
\end{equation}
Here $z=\psi_0/\psi_1$ is the ratio of the two amplitudes: the
decoded quantity in Eq.~\eqref{eq:decode} is a \emph{ratio}, not an
amplitude, and it ranges over all of $\CP$ even though the state
stays normalized.  The map sends $z=0$ to the south pole $\ket{1}$,
$z=\infty$ to the north pole $\ket{0}$, and the equator to
$|z|=1$; its Bloch vector is
$(\langle\PauliX\rangle,\langle\PauliY\rangle,\langle\PauliZ\rangle)
=(2\re z,-2\im z,|z|^2-1)/(1+|z|^2)$.

Because decoding reads a ratio, a single-qubit gate acts on the
encoded point by a M\"obius transformation: a unitary
$\bigl(\begin{smallmatrix}a&b\\c&d\end{smallmatrix}\bigr)$ sends
$\ket{z}$ to $\ket{(az+b)/(cz+d)}$.  The Pauli gates give
$\PauliX\colon z\mapsto 1/z$, $\PauliY\colon z\mapsto -1/z$, and
$\PauliZ\colon z\mapsto -z$, and the phase gate $e^{i\phi\PauliZ}$
acts as $z\mapsto e^{2i\phi}z$.  The state $\ket{z}$ is prepared
from $\ket{0}$ by
\begin{equation}
	U_z=\frac{1}{\sqrt{1+|z|^2}}
	\begin{pmatrix} z & 1 \\ 1 & -\bar z\end{pmatrix},
	\label{eq:prep-unitary}
\end{equation}
whose columns are $\ket{z}$ and the antipodal state
$\ket{-1/\bar z}$.

\subsection{The signal operator and its base case}
\label{subsec:signal}

A QSP protocol interleaves a signal-dependent gate with tunable
phases.  In the stereographic chart the signal is a point
$z=re^{i\theta}\in\CP$, and we take as signal operator the
unnormalized matrix
\begin{equation}
	O_z=\begin{pmatrix} z & -1 \\ 1 & \bar z\end{pmatrix},
	\qquad \det O_z=1+|z|^2,
	\label{eq:Oz}
\end{equation}
whose determinant is the squared Fubini--Study norm of the
homogeneous coordinate $(z:1)$.  For a real signal $z=x\in\R$ the
operator $O_x=\bigl(\begin{smallmatrix}x&-1\\1&x\end{smallmatrix}\bigr)$
is $\sqrt{1+x^2}$ times a rotation by $\theta=\mathrm{arccot}\,x$, and
its $k$th power has entries
\begin{equation}
	O_x^{\,k}\ \propto\
	\begin{pmatrix}\cos k\theta & -\sin k\theta\\
		\sin k\theta & \cos k\theta\end{pmatrix},
	\qquad \theta=\mathrm{arccot}\,x,
	\label{eq:Ox-power}
\end{equation}
after dividing by $(1+x^2)^{k/2}$.  The matrix elements are the
\emph{rational Chebyshev functions} of Boyd~\cite{Boyd1990}
($\TB_k$ on the diagonal, $\mp\SB_k$ off it),
\begin{equation}
	\TB_k(x)=\cos(k\,\mathrm{arccot}\,x),\qquad
	\SB_k(x)=\sin(k\,\mathrm{arccot}\,x),
	\label{eq:rat-cheb}
\end{equation}
which are the images of $\cos k\theta$ and $\sin k\theta$ under the
projection $x=\cot\theta$ and are orthogonal on $\R$ with the
Cauchy weight $1/(1+x^2)$~\cite{Boyd1990,Boyd1987}.  Throughout we
fix the branch $\mathrm{arccot}\colon\R\to(0,\pi)$, so $\sin\theta>0$
and $\sqrt{1+x^2}=1/\sin\theta$ carries no sign ambiguity.  Conjugating
the normalized $k$th power by $V=\mathrm{diag}(1,i)$, that is forming
$V\,[(1+x^2)^{-k/2}O_x^{\,k}]\,V^{-1}$, gives the standard QSP signal
$W(\tilde x)^k$ of Remark~\ref{rem:basis-change}, with $+i\sin k\theta$
in \emph{both} off-diagonal entries.  Boyd writes
$\TB_k$ as $T_k$ composed with the compression
$x\mapsto x/\sqrt{1+x^2}$, since $T_k(\cos\theta)=\cos k\theta$;
Eq.~\eqref{eq:rat-cheb} is the same family.

\subsection{Completeness in the stereographic chart}
\label{subsec:completeness}

The signal operator $O_x$ and the phase gates generate a
QSP sequence whose achievable functions we now characterize.
Throughout, $P^*$ denotes the polynomial
$P^*(x)=\overline{P(\bar x)}$ obtained by conjugating the
coefficients of $P$, so that $P^*$ coincides with the complex
conjugate $\bar P$ on real arguments.  We give the real-signal case
first.

\begin{theorem}[Real stereographic QSP]
	\label{thm:real-qsp}
	For $x\in\R$ let
	$U_{\bm\phi}(x)=e^{i\phi_0\PauliZ}\prod_{j=1}^{d}O_x\,
	e^{i\phi_j\PauliZ}$.  A pair of polynomials $P,Q\in\C[x]$ is
	realizable as
	\begin{equation}
		U_{\bm\phi}(x)=
		\begin{pmatrix} P(x) & -Q^*(x)\\ Q(x) & P^*(x)\end{pmatrix}
		\label{eq:real-qsp-form}
	\end{equation}
	for some $\bm\phi\in\R^{d+1}$ if and only if
	\begin{enumerate}
		\item[(i)] $|P(x)|^2+|Q(x)|^2=(1+x^2)^d$ for all $x\in\R$;
		\item[(ii)] $P$ has parity $d\bmod 2$ and $Q$ has parity
			$(d-1)\bmod 2$;
		\item[(iii)] $\deg P\le d$ and $\deg Q\le d-1$, with the
			convention $\deg Q=-1\Leftrightarrow Q\equiv 0$.
	\end{enumerate}
\end{theorem}

The forward direction is an induction on $d$: each factor
$O_x\,e^{i\phi\PauliZ}$ raises the degree by one and multiplies the
normalization by $\det O_x=1+x^2$, which reproduces (i)--(iii).
The backward direction is the substantive half.  Condition (i) forces
$|P(i)|^2=|Q(i)|^2$ by evaluation at $x=i$; choosing the phase
$e^{2i\phi}=-iP(i)/Q^*(i)$ makes the reduced polynomials
$\widetilde P,\widetilde Q$ vanish at $x=i$, and, by parity, at
$x=-i$.  Division by $1+x^2$ then lowers the degree by one and the
induction closes.  Appendix~\ref{app:thm-real} gives the details.
The normalization (i) is the only change from the standard QSP
theorem: it is $(1+x^2)^d$ rather than $1$, because the signal
operator $O_x$ is unnormalized.  The physical gate is the unitary
$(1+x^2)^{-d/2}U_{\bm\phi}(x)\in\SU(2)$, and this global factor
cancels in the decoded ratio $P/Q$; after the basis change of
Remark~\ref{rem:basis-change} the constraint (i) becomes the
standard completeness constraint $|P|^2+|Q|^2=1$ in the compressed
variable.

Lifting the controls from $\PauliZ$-phases to arbitrary $\SU(2)$
rotations, and the signal from $\R$ to $\CP$, gives the generalized
theorem.  Its achievable class consists of \emph{bivariate}
polynomials in $z$ and $\bar z$.

\begin{theorem}[Generalized stereographic QSP]
	\label{thm:gen-qsp}
	Let $O_z$ be as in Eq.~\eqref{eq:Oz} and let
	$G_0,\dots,G_d\in\SU(2)$.  A pair $P,Q\in\C[z,\bar z]$, where
	now $P^*(z,\bar z)=\overline{P(\bar z,z)}$ conjugates coefficients
	and swaps $z\leftrightarrow\bar z$, is realizable as
	\begin{equation}
		G_d\,O_z\,G_{d-1}\,O_z\cdots O_z\,G_0
		=\begin{pmatrix} P(z,\bar z) & -Q^*(z,\bar z)\\
			Q(z,\bar z) & P^*(z,\bar z)\end{pmatrix}
		\label{eq:gen-qsp-form}
	\end{equation}
	if and only if
	\begin{enumerate}
		\item[(i)] every monomial $z^j\bar z^k$ of $P$ and $Q$ has
			$j+k\le d$;
		\item[(ii)] $P P^*+Q Q^*=(1+|z|^2)^d$.
	\end{enumerate}
\end{theorem}

The forward direction is again an induction on the number of
factors.  The backward direction reads the leading data off a
target pair and strips one layer at a time.  Applying the inverse
iterate $O_z^{-1}G_d^{-1}$ produces reduced entries
$P'=z^*\widetilde P+\widetilde Q$ and $Q'=z\widetilde Q-\widetilde
P$; requiring $1+z\bar z$ to divide both fixes the control $G_d$,
and the quotient lowers the bidegree from $d$ to $d-1$.  The
divisibility rests on a factorization in the Laurent ring:

\begin{lemma}
	\label{lem:DDstar}
	If $D\in\C[z,z^{-1}]$ satisfies $D D^*=-1$, then $D=c\,z^k$ with
	$|c|=1$ and $k$ odd.
\end{lemma}

The lemma pins down the units of the ring that can appear, and with
it the case analysis of the backward direction closes; the full
argument, including the unique-factorization step, is in
Appendix~\ref{app:thm-gen}.  Theorem~\ref{thm:gen-qsp} is the
stereographic form of generalized QSP~\cite{MotlaghWiebe2023}: the
bivariate structure in $(z,\bar z)$ is the projective-coordinate
image of the Laurent structure in $z=e^{i\theta}$, and the
Fubini--Study normalization $(1+|z|^2)^d$ is the image of the
constraint $|P|^2+|Q|^2=1$ on the unit circle.

\subsection{Three readings of the same circuit}
\label{subsec:remarks}

\begin{remark}[Reduction to standard QSP]
	\label{rem:basis-change}
	The real signal operator, once normalized, is the standard
	QSP signal conjugated by a fixed gate.  With
	$V=\mathrm{diag}(1,i)$, $\tilde x=x/\sqrt{1+x^2}$, and
	$W(\tilde x)=\bigl(\begin{smallmatrix}\tilde x &
	i\sqrt{1-\tilde x^2}\\ i\sqrt{1-\tilde x^2}&\tilde
	x\end{smallmatrix}\bigr)$ the standard signal,
	\begin{equation}
		(1+x^2)^{-1/2}O_x=V^{-1}W(\tilde x)V,
	\end{equation}
	equivalently $W(\tilde x)=V\,[(1+x^2)^{-1/2}O_x]\,V^{-1}$, and $V$
	commutes with every phase gate.  A stereographic QSP
	circuit is therefore a standard QSP circuit conjugated by a
	gate that depends on neither the signal nor the phases.  Two
	consequences follow.  The diagonal matrix element is unchanged,
	so the completeness theorem and every phase-finding
	algorithm~\cite{DongMengWhaleyEtAl2021,NiSarkarYingLin2025}
	transfer verbatim.  And the unboundedness of the decoded ratio
	is produced by the decoding step of Eq.~\eqref{eq:decode}, not
	by the circuit, whose matrix elements never exceed one.
\end{remark}

\begin{remark}[A hardware-executable read-out]
	\label{rem:measurement}
	After the basis change of Remark~\ref{rem:basis-change} the
	circuit acts on real amplitudes, so its output is a
	real-amplitude state, with a single flag qubit carrying the
	imaginary unit in the real-number formulation of
	Ref.~\cite{BarriosHita2026}.  Such states are read out by the
	$n+1$-setting protocol of
	Ref.~\cite{NguemtoLeytonOrtega2026}: computational-basis
	sampling and one Hadamard-rotated setting per qubit, which is
	minimal.  The stereographic ratio is then formed
	classically.  The read-out costs $n+1$ measurement settings and
	$\mathcal{O}(r^4/\varepsilon^2)$ shots to recover the decoded
	value at signal magnitude $r$ to absolute precision
	$\varepsilon$, the inverse projection amplifying sampling error
	by $(1+r^2)/2$ (Appendix~\ref{app:readout}).  This makes the
	decoding interface executable for order-one signals; it is a
	trade of circuit depth for shot count, not an asymptotic
	speedup, and its systematic use is left to future work.
\end{remark}

\begin{remark}[One geometry, two signal operators]
	\label{rem:two-signals}
	The stereographic map appears twice in this paper.  At the
	single-qubit scale it selects the signal operator $O_z$, whose
	powers produce the rational Chebyshev functions and whose
	achievable class the two theorems above characterize.  At the
	multi-qubit scale (Sec.~\ref{sec:stateprep}) the same map sends
	a uniform angular grid $\theta_j$ to the cotangent grid
	$y_j=L\cot\theta_j$ on $\R$, on which a diagonal walk operator
	$e^{iA}=\sum_j e^{i\theta_j}\ket{j}\!\bra{j}$ acts and the
	rational Chebyshev basis pulls back to the Fourier basis.  The
	two operators $O_z$ and $e^{iA}$ are the \emph{same geometry},
	not the same gate: they are the single- and multi-qubit
	incarnations of one stereographic construction, and the
	pairing of chart, grid, and signal operator is what turns the
	approximation theory of this section into the gate counts of
	Sec.~\ref{sec:stateprep}.
\end{remark}
